\section{Compact clips leave half the mass outside}\label{sec:clips}

In this section, $E$ is an infinite-dimensional closed real subspace of
a probability $L^1$ space, with its inherited norm. The clipping
argument uses only total boundedness of $B_E$ in measure.

For $x\in S_E$, $y\in E$, and $t>0$, the real scalar identity
$(|a+b|+|a-b|)/2=\max\{|a|,|b|\}$ gives
\begin{equation}\label{eq:overlap-identity}
 \frac{\|x+ty\|_1+\|x-ty\|_1}{2}-1
 =\int(t|y|-|x|)_+\,d\mathbb P
 =t\|y\|_1-\int\min\{|x|,t|y|\}\,d\mathbb P.
\end{equation}
Here $r_+=\max\{r,0\}$.  The lower bound therefore amounts to controlling
the overlap with $x$.

\begin{lemma}\label{lem:clips}
Let $E$ be an infinite-dimensional closed real subspace of
$L^1(\Omega,\mathbb P)$, where $\mathbb P$ is a probability measure.
If its unit ball $B_E$ is totally bounded in measure, then
$H_E(x,t)\ge\max\{t/2,t-1\}$ for every $x\in S_E$ and $t>0$.
\end{lemma}
\begin{proof}
Fix $x$ and $t$.  Clip a function $y$ at the pointwise height $|x|/t$:
\[
 C(y)=\operatorname{sgn}(y)\min\{|y|,|x|/t\}.
\]
The clipped functions need not belong to $E$.  Their norm closure
$K=\overline{C(B_E)}^{\,L^1}$ is compact.  Indeed, every sequence in $B_E$
has a subsequence Cauchy in measure.  Clipping is pointwise
$1$-Lipschitz, so the clipped subsequence is also Cauchy in measure.
For two of its terms $b_i,b_j$ and any $a>0$,
\[
 \|b_i-b_j\|_1
 \le a+\int_{\{|b_i-b_j|>a\}}2|x|/t\,d\mathbb P.
\]
Choose $a$ small, then use absolute continuity of the integral and the
measure-Cauchy property.  This proves that the subsequence is Cauchy in
$L^1$.  Completeness proves the compactness assertion.

Fix $\varepsilon>0$ and take a finite $\varepsilon$-net
$b_1,\ldots,b_m$ of $K$.  Set $\sigma_i=\operatorname{sgn}(b_i)$ and
\[
 F=\bigcap_{i=1}^m\ker\left(y\longmapsto
                 \int\sigma_i y\,d\mathbb P\right)\subseteq E.
\]
This is a closed subspace of finite codimension.  For $y\in S_F$, choose
$i$ with $\|C(y)-b_i\|_1<\varepsilon$.  Clipping preserves signs and
reduces absolute values, so
$\|y-C(y)\|_1=1-\|C(y)\|_1$.  Since $\int\sigma_i y=0$, we obtain
\[
 \|C(y)\|_1-\varepsilon
 \le\|b_i\|_1
 =\left|\int\sigma_i(b_i-y)\,d\mathbb P\right|
 \le1-\|C(y)\|_1+\varepsilon.
\]
Thus $\|C(y)\|_1\le1/2+\varepsilon$.
Equation~\eqref{eq:overlap-identity} now gives averaged increment at
least $t/2-t\varepsilon$ on $S_F$.  The same bound holds for
$\|y\|_1\ge1$ by radial monotonicity.  Taking the supremum over $F$ and
letting $\varepsilon\downarrow0$ gives $H_E(x,t)\ge t/2$.
The subspace is allowed to change with $\varepsilon$; this step does not
assert that one subspace attains the limiting bound.

For every $\|y\|_1\ge1$, the triangle inequality also gives
\[
 \frac{\|x+ty\|_1+\|x-ty\|_1}{2}-1\ge t\|y\|_1-1\ge t-1.
\]
Combining the two bounds proves the lemma.
\end{proof}

\begin{remark}[A separated-sequence proof]\label{rem:separated}
There is also a short contradiction argument giving the same half-mass
bound. Fix a unit center $x$, a radius $t>0$, and $0<a<1/2$.
Suppose no subspace of finite codimension had averaged increment at
least $at$ in every direction of norm at least one. By radial monotonicity and
Equation~\eqref{eq:overlap-identity}, every subspace of finite
codimension would then contain a unit vector $y$ with
$\|y-C(y)\|_1<a$. Inductively choose such $y_j$ in the kernels of
norm-one functionals norming the preceding vectors.  Then
$\|y_i-y_j\|_1\ge1$ for $i\ne j$.  Compactness of the clipped set gives
a subsequence whose clips are Cauchy in norm, whereas along this
subsequence
\[
 \|y_i-y_j\|_1
 \le\|y_i-C(y_i)\|_1+\|C(y_i)-C(y_j)\|_1
          +\|C(y_j)-y_j\|_1<2a+o(1)<1,
\]
a contradiction for sufficiently late pairs. Thus this argument yields
$H_E(x,t)\ge t/2$ by taking the supremum over the corresponding subspaces
and then letting $a\uparrow1/2$. The finite-net proof above identifies
an explicit finite family of functionals that removes the overlap.
\end{remark}

\section{Almost diametral points in weak neighborhoods}\label{sec:weak-neighborhoods}

To find almost diametral points subject to finitely many linear
constraints, we pass from slices to relatively weakly open subsets of
the ball. These points will supply both the upper bounds for the moduli
and the obstruction to AUC renorming. A \emph{slice} of $B_X$ is a nonempty set
\[
 H=\{y\in B_X:\ell(y)>\alpha\},\qquad
 \ell\in X^*,\quad\alpha\in\mathbb R.
\]
We first prove the geometric characterization of the Daugavet property
due to Kadets, Shvidkoy, Sirotkin, and Werner
\cite[Lemma~2.2]{KSSW2000}, also recalled by Kadets and Werner
\cite[Lemma~1.1]{KW2004}.

\begin{lemma}[Daugavet slice characterization]\label{lem:daugavet-slices}
A nonzero real Banach space $X$ has the Daugavet property if and only if,
for every $w\in S_X$, every $\rho>0$, and every slice $H$ of $B_X$,
there is $y\in H$ with $\|w+y\|\ge2-\rho$.
\end{lemma}
\begin{proof}
Suppose first that $X$ has the Daugavet property. If the functional
defining $H$ is zero, then $H=B_X$ and we may take $y=w$. Otherwise,
after dividing the functional and threshold by its norm, write
$H=\{y\in B_X:\ell(y)>\alpha\}$ with $\|\ell\|=1$ and $\alpha<1$.
Take $0<\eta<\min\{1-\alpha,\rho/2\}$. The rank-one operator
$Ty=\ell(y)w$ has norm one, so there is $y\in S_X$ with
$\|y+\ell(y)w\|>2-\eta$. Replacing $y$ by $-y$ if necessary makes
$\ell(y)\ge0$ without changing this inequality. It follows that
$\ell(y)>1-\eta>\alpha$ and
\[
 \|y+w\|\ge\|y+\ell(y)w\|-(1-\ell(y))>2-2\eta>2-\rho.
\]

Conversely, assume the slice condition and let
$Tz=\ell(z)v$ be a nonzero rank-one operator with $\|\ell\|=1$.
Set $a=\|v\|=\|T\|$ and $w=v/a$. For $0<\eta<1$, the slice condition
provides $y\in B_X$ with $\ell(y)>1-\eta$ and
$\|w+y\|\ge2-\eta$. A norm-one functional $\varphi$ norming $w+y$
takes value at least $1-\eta$ on each summand. Consequently,
\[
 \|I+T\|\ge\|y+Ty\|
 \ge\varphi(y)+a\ell(y)\varphi(w)
 \ge(1-\eta)+a(1-\eta)^2.
\]
Let $\eta\downarrow0$ and use $\|I+T\|\le1+\|T\|$.
The identity is immediate when $T=0$.
\end{proof}

The passage from slices to weak neighborhoods uses Bourgain's lemma
on convex combinations of slices
(see \cite{ANP2019} for the attribution).
We give a finite-dimensional proof.  Its use in passing to weak
neighborhoods follows the same pattern as Shvydkoy's proof of the
weak-open characterization \cite[Lemma~3]{Shvydkoy2000}.

\begin{lemma}[Bourgain's lemma]\label{lem:slice-combination}
Every nonempty relatively weakly open subset $U$ of a real Banach
space's closed unit ball contains a convex combination of nonempty
slices of that ball.
\end{lemma}
\begin{proof}
Choose $u\in U$, a continuous linear map $T:X\to\mathbb R^k$, and an
open set $O\subseteq\mathbb R^k$ such that
$u\in B_X\cap T^{-1}(O)\subseteq U$.  Put $D=\overline{T(B_X)}$.  We want to place a convex combination of
small slices of $D$ inside $O$ and then pull them back to $B_X$.
For this it suffices to approximate $T(u)$ by convex combinations of
exposed points, that is, unique maximizers of linear functionals on
$D$.  We first verify that $D$ is the closed convex hull of those points.
Put $M=\max_{q\in D}|q|$, using a Euclidean norm.  For a unit vector $a$
and $R>M$, choose $p_R\in D$ farthest from $-Ra$.  Expanding the squares
in $|q+Ra|\le|p_R+Ra|$ gives
\[
 2(p_R+Ra)\cdot(q-p_R)\le-|q-p_R|^2\qquad(q\in D).
\]
Thus $p_R$ is the unique maximizer of the linear functional with
coefficient $p_R+Ra$.  If $q_a$ maximizes $a\cdot q$ on $D$, the same
comparison gives
\[
 0\le a\cdot q_a-a\cdot p_R
 \le\frac{|p_R|^2-|q_a|^2}{2R}\le\frac{M^2}{2R}.
\]
The exposed points therefore have the same supremum as $D$ for every
linear functional.  Separation proves the closed-convex-hull assertion,
also when $D$ is a singleton.

Since $T(u)\in O$, there are exposed points $p_1,\ldots,p_l$ and
positive weights $\lambda_i$ summing to one such that
$p=\sum_i\lambda_i p_i\in O$.  Choose $r>0$ with the Euclidean ball
$B(p,r)\subseteq O$.  If $a_i\cdot q$ exposes $p_i$, compactness and
uniqueness of its maximum give a threshold $\beta_i<a_i\cdot p_i$ with
\[
 \{q\in D:a_i\cdot q>\beta_i\}\subseteq B(p_i,r).
\]
The sets $H_i=\{y\in B_X:a_i\cdot T(y)>\beta_i\}$ are nonempty slices,
since $T(B_X)$ is dense in $D$.  For $y_i\in H_i$,
\[
 \left|T\left(\sum_i\lambda_i y_i\right)-p\right|
 \le\sum_i\lambda_i|T(y_i)-p_i|<r.
\]
Their convex combination consequently belongs to $U$.
\end{proof}

The slice property can now be applied successively within the slices
provided by Bourgain's lemma.  This gives the following specialization
of Shvydkoy's weak-open result \cite[Lemma~3]{Shvydkoy2000}.

\begin{lemma}\label{lem:diametral}
If $X$ has the Daugavet slice property, then, for every $x\in S_X$ and
every relatively weakly open neighborhood $U$ of $x$ in $B_X$,
\[
 \sup_{v\in U}\|v-x\|=2.
\]
\end{lemma}
\begin{proof}
By Lemma~\ref{lem:slice-combination}, choose slices $H_i$ and positive
weights $\lambda_i$ such that $\sum_i\lambda_i=1$ and
$\sum_i\lambda_iH_i\subseteq U$.  Fix $0<\eta<1$ and set $z_0=-x$.
At step $i$, apply the slice property at $z_{i-1}/\|z_{i-1}\|$ to choose
$y_i\in H_i$ with
$\|z_{i-1}/\|z_{i-1}\|+y_i\|\ge2-\eta$.
A norm-one functional norming this sum takes value at least $1-\eta$
on each of its two summands.  Therefore, with
$z_i=z_{i-1}+\lambda_i y_i$,
\[
 \|z_i\|\ge(1-\eta)(\|z_{i-1}\|+\lambda_i)>0.
\]
The induction is well defined and yields
\[
 \|z_l\|\ge(1-\eta)^l+
    \sum_{i=1}^l\lambda_i(1-\eta)^{l-i+1}
 \ge2(1-\eta)^l.
\]
Now $v=\sum_i\lambda_i y_i\in U$ and $z_l=v-x$.  Let
$\eta\downarrow0$; the reverse bound follows from $v,x\in B_X$.
\end{proof}

\section{One family of directions gives both curves}\label{sec:curves}

\begin{proof}[Proof of the formulas in Theorem~\ref{thm:exact}]
Lemma~\ref{lem:clips} gives the averaged lower bound.  For the one-sided
bound, choose a norm-one functional $\varphi$ with
$\varphi(x)=1$.  Every $y\in S_{\ker\varphi}$ satisfies
\[
 \|x+ty\|\ge1,\qquad \|x+ty\|\ge t-1.
\]
Hence $D_E(x,t)\ge\max\{0,t-2\}$.

For the upper bounds, fix any $F\in\operatorname{cof}(E)$ and choose a
bounded finite-rank projection $P$ with kernel $F$.  Such a projection
is obtained by composing the quotient map with a linear section of
$E/F$.  For $0<\varepsilon<1/2$,
\[
 U_\varepsilon=\{v\in B_E:\|P(v-x)\|<\varepsilon\}
\]
is a relatively weakly open neighborhood of $x$.  The finite-dimensional
range of $P$ makes it weak-to-norm continuous.  By
Lemma~\ref{lem:diametral}, choose $v_\varepsilon\in U_\varepsilon$ with
$\|v_\varepsilon-x\|>2-\varepsilon$, and put
\[
 z_\varepsilon=(I-P)(v_\varepsilon-x),\quad
 L_\varepsilon=\|z_\varepsilon\|,\quad
 y_\varepsilon=z_\varepsilon/L_\varepsilon\in S_F.
\]
Then
\[
 2-2\varepsilon<L_\varepsilon\le2+\varepsilon,
 \qquad \|x+z_\varepsilon\|\le1+\varepsilon.
\]
Thus the directions lie exactly in $F$, their unnormalized lengths tend
to two, and the translated endpoints remain close to $B_E$.

For any $t>0$, set $\lambda=\min\{t/L_\varepsilon,1\}$.  The identity
\[
 x+t y_\varepsilon
 =(1-\lambda)x+\lambda(x+z_\varepsilon)
       +(t-L_\varepsilon)_+y_\varepsilon
\]
and the triangle inequality give the single bound
\[
 \|x+t y_\varepsilon\|
 \le1+\varepsilon+(t-L_\varepsilon)_+.
\]
The right-hand side converges to $1+(t-2)_+=\max\{1,t-1\}$.  Since
$\|x-t y_\varepsilon\|\le1+t$, the one-sided and averaged increments
have upper limits at most
\[
 (t-2)_+,\qquad \frac{t+(t-2)_+}{2},
\]
respectively.  These are $\max\{0,t-2\}$ and
$\max\{t/2,t-1\}$.  In particular the argument includes $t=2$, where
the two branches meet.  Each inner infimum for the fixed $F$ is at most
the corresponding limit.  Taking the supremum over $F$ proves the
upper bounds.  Both formulas hold for every unit center, so taking the
infimum over centers leaves them unchanged.
\end{proof}

\section{Why a different norm cannot be AUC}\label{sec:renorming}

The one-sided modulus vanishes for $0<t\le2$, so the inherited norm is not AUC.
Ruling out all equivalent norms requires the weak-neighborhood
property itself. The following argument applies to any infinite-dimensional
real Banach space in which every relatively weakly open neighborhood of a unit
vector contains points of the unit ball arbitrarily close to distance
two from that vector, as in Lemma~\ref{lem:diametral}.

\begin{proof}[Proof of the renorming assertion in Theorem~\ref{thm:exact}]
Suppose $N$ is an equivalent AUC norm.  Choose $a>0$ with
$N(u)\ge a\|u\|$ for all $u$, and put
$R=\sup_{u\in B_E}N(u)<\infty$.  Fix, in this order,
\[
 t_0=\frac{a}{4R},\qquad
 0<\gamma<\overline\delta_N(t_0),\qquad
 x_0\in S_E\ \text{with}\ (1+\gamma)N(x_0)>R.
\]
The last choice is possible since $R=\sup_{u\in S_E}N(u)$.
The definition of $\overline\delta_N(t_0)$ at $x_0/N(x_0)$ supplies
$F\in\operatorname{cof}(E)$ such that
\begin{equation}\label{eq:auc-outside}
 N(x_0+z)\ge(1+\gamma)N(x_0)
 \quad\text{if }z\in F,\quad N(z)\ge t_0N(x_0).
\end{equation}
To justify the entire range of radii, convexity gives, for $N(w)=1$ and
$r\ge t_0$,
\[
 N(x_0/N(x_0)+rw)-1
 \ge\frac r{t_0}\bigl(N(x_0/N(x_0)+t_0w)-1\bigr).
\]

Choose
$0<\eta<\min\{(1+\gamma)N(x_0)-R,a/4\}$.  The set
\[
 U=\{u\in B_E:\inf_{z\in F}N(u-x_0-z)<\eta\}
\]
is a relatively weakly open neighborhood of $x_0$: the infimum is the
norm in the finite-dimensional quotient $E/F$.  Equivalent norms have
the same continuous functionals, hence the same weak topology.
For $u\in U$, choose $z\in F$ with $N(u-x_0-z)<\eta$.
If $N(z)\ge t_0N(x_0)$, Equation~\eqref{eq:auc-outside} gives
\[
 N(u)>(1+\gamma)N(x_0)-\eta>R,
\]
contrary to $u\in B_E$.  Thus
\[
 N(u-x_0)<t_0N(x_0)+\eta\le t_0R+\eta<a/2,
 \qquad\|u-x_0\|<1/2.
\]
This contradicts Lemma~\ref{lem:diametral}.
\end{proof}
